% Section 17: 01-full-contraction | Part: 3 | Depends on: see docs/REWRITE_PLAN.md | Status: todo | Original lines: 2613-2635
%!TEX root = ../../../main.tex
Let $\phi:U\subset\R^k\to M$ be a parametrised $k$-dimensional piece of $M$, and $\alpha$ a $k$-form. Then
\begin{equation}\label{eq:intk}
\int_\phi\alpha:=\int_U\phi\pb\alpha=\int_U\alpha_{\phi(u)}\bigl(\partial_1\phi,\dots,\partial_k\phi\bigr)\,\dd u^1\cdots\dd u^k,
\end{equation}
that is, evaluate $\alpha$ on the tangent $k$-frame of the parametrisation (a \emph{full contraction}) and integrate the resulting function with the Lebesgue measure. For $k=1$ this is \eqref{eq:lineint}.
The same reasoning as in Proposition~\ref{prop:reparam} shows invariance under orientation-preserving reparametrisations (the Jacobian determinant appears through the multilinear \emph{alternating} character of $\alpha$), which is why the integrands of oriented $k$-dimensional objects must be alternating $k$-forms. On an oriented manifold one globalises with a partition of unity.

\emph{Why a determinant appears.} For a $2$-form $\alpha=a\,\dd x\wedge\dd y$ and a parametrisation $\phi(u,v)=(x(u,v),y(u,v))$, the tangent frame is $\partial_u\phi=(x_u,y_u)$, $\partial_v\phi=(x_v,y_v)$, so
\[
\alpha(\partial_u\phi,\partial_v\phi)=a\,(x_uy_v-y_ux_v)=a\det D\phi,
\]
where $D\phi$ is the Jacobian matrix. Therefore \eqref{eq:intk} is the change-of-variables formula $\int_\phi a\,\dd x\wedge\dd y=\iint a(\phi(u,v))\det D\phi\,\dd u\,\dd v$ with the \emph{sign} of the Jacobian kept: orientation-preserving reparametrisations ($\det>0$) leave it unchanged and orientation-reversing ones flip its sign.

\begin{example}[A surface integral, step by step]\label{ex:paraboloid}
Integrate $\alpha=z\,\dd x\wedge\dd y$ over the paraboloid cap $\phi(u,v)=(u,v,u^2+v^2)$, $u^2+v^2\le1$.
\begin{enumerate}
\item Pull back the coordinate functions: $\phi\pb x=u$, $\phi\pb y=v$, $\phi\pb z=u^2+v^2$; since pull-back commutes with $\dd$ and $\wedge$, $\phi\pb\dd x=\dd u$ and $\phi\pb\dd y=\dd v$.
\item Hence $\phi\pb\alpha=(u^2+v^2)\,\dd u\wedge\dd v$.
\item Evaluate on the frame $(\partial_u,\partial_v)$: $(\dd u\wedge\dd v)(\partial_u,\partial_v)=1$, so $\int_\phi\alpha=\iint_{u^2+v^2\le1}(u^2+v^2)\,\dd u\,\dd v$.
\item Polar coordinates (the factor $r$ is the determinant of Example~\ref{ex:polarwedge}): $\displaystyle\int_0^{2\pi}\!\!\int_0^1r^2\cdot r\,\dd r\,\dd\theta=2\pi\cdot\tfrac14=\tfrac\pi2$.
\end{enumerate}
\end{example}
